\documentclass{article}

\usepackage{booktabs}
\usepackage{tabularx}
\usepackage{array}
\usepackage{float}
\usepackage{enumitem}

\title{Stakeholder Interview Report\\\progname}

\author{\authname}

\date{}

\input{../../Comments.text}
\input{../../Common.text}

\newcolumntype{L}[1]{>{\raggedright\arraybackslash}p{#1}}
\newcolumntype{Y}{>{\raggedright\arraybackslash}X}
\setlength{\tabcolsep}{4pt}

\begin{document}

\maketitle

\begin{table}[hp]
\caption{Revision History} \label{TblRevisionHistory}
\begin{tabularx}{\textwidth}{llX}
\toprule
\textbf{Date} & \textbf{Developer(s)} & \textbf{Change}\\
\midrule
11/10/2026 & Pedro Azevedo & Template created for report \\
\bottomrule
\end{tabularx}
\end{table}

\newpage
\tableofcontents
\listoftables
\newpage

% =====================================================================
\section{Introduction}
\label{sec:intro}

\subsection{Purpose}

The purpose of the interviews was to collect specified and verifiably valid feedback from the 
users with the most experience in the club. With the base of the project identified, outlined 
in the SRS, the interviews provided proper guiding information about the prioritization and 
further specification of services. The prior elicitation provides data to contrast against the 
interview, ensuring that experience member results are not overly biased. Product information 
that is related or benefits from input of experienced members, or executive specific 
requirements will be addressed through the interviews. There will also be discussion 
regarding the available data for testing, specifically data that has more records or 
background stored from current experience members.

\subsection{Relation to Other Documents}
\todo{Member survey and SRS P.7 \cite{TroitskySRS}; VnV Plan Custom Documents
section \cite{TroitskyVnVPlan}; Winter usability report.}

% =====================================================================
\section{Methodology}
\label{sec:method}

\subsection{Interview Format}
The Fall CD interviews will be semi-structured, with questions tailored
primarily to experienced members, and additionally to specific roles.
Interviews will be around 10--15 minutes long, and conducted with 2--3
participants at a time, allowing for both discussion and individual responses.

The interview guide is intended to be flexible allowing for experienced members
to direct the discussion to areas they feel are important, while still ensuring
that core concerns are addressed. As many of the interviewed members will be a
part of many of these user groups, certain interviews will cover multiple
question sets.

\subsection{Participant Selection and Recruitment}
\todo{Stakeholder groups targeted, selection criteria, scheduling.}

\subsection{Session Procedure and Roles}
\todo{Interviewer and notetaker roles, introduction script, how notes were
recorded (interview record template).}

\subsection{Ethics and Anonymization}
\todo{Consent, voluntary participation, what was and was not recorded, how
transcripts were anonymized.}

\subsection{Analysis Method}
\todo{Empathy map per stakeholder group; thematic coding against SRS sections.}

% =====================================================================
\section{Interview Questions}
\label{sec:questions}

\begin{table}[H]
\centering
\caption{Interview questions by stakeholder group}
\label{tab:questions}
\small
\begin{tabularx}{\textwidth}{L{0.9cm} Y L{2.6cm}}
\toprule
\textbf{ID} & \textbf{Question} & \textbf{Traces to} \\
\midrule
\multicolumn{3}{l}{\textit{All participants, general questions}} \\
G1 & What is your role in the club, and how many years have you been involved? & \\
G2 & Which tools do your teams use to design, analyze, and store their bridges? & \\
G3 & What were the largest challenges you faced when you first joined the club? & \\
\midrule
\multicolumn{3}{l}{\textit{Executives}} \\
E1 & What information regarding the club administration is most important to consolidate? & \\
E2 & When you first became an executive, what were your largest difficulties? & \\
E3 & How would you like to see the software sub team integrated into the club administration? & \\
\midrule
\multicolumn{3}{l}{\textit{Team captains and experienced members}} \\
C1 & What tool have you found is the most accurate and useful in predicting your bridge's performance? & \\
C2 & How close to a model estimate do you expect your bridge's performance to be, including uncertainties? & \textsc{failure\_load\_\allowbreak tolerance} \\
C3 & What are the most common reasons for bridge failures in your experience? & \\
C4 & Does the club have models and recorded failure loads of past bridges that could be used to test the Toolbox? & \\
\midrule
\multicolumn{3}{l}{\textit{Closing, all participants}} \\
CL1 & Is there anything about the club's information or tools that we did not ask about? & \\
CL2 & Are there any items that you feel are important to include in the Toolbox that we have not discussed? & \\
\bottomrule
\end{tabularx}
\end{table}

% =====================================================================
\section{Participants}
\label{sec:participants}

\begin{table}[H]
\centering
\caption{Participants (non-identifying)}
\label{tab:participants}
\small
\begin{tabularx}{\textwidth}{L{0.8cm} L{1.3cm} Y L{1.6cm} L{2.6cm}}
\toprule
\textbf{ID} & \textbf{Session} & \textbf{Role(s)} & \textbf{Years in club} & \textbf{Modelling experience} \\
\midrule
P1 & & & & \\
P2 & & & & \\
P3 & & & & \\
\bottomrule
\end{tabularx}
\end{table}

% =====================================================================
\section{Interview Sessions}
\label{sec:sessions}

\begin{table}[H]
\centering
\caption{Interview sessions}
\label{tab:sessions}
\small
\begin{tabularx}{\textwidth}{L{1.3cm} L{1.4cm} L{2.1cm} L{2.0cm} L{2.0cm} Y}
\toprule
\textbf{Session} & \textbf{Date} & \textbf{Participants} & \textbf{Interviewer} & \textbf{Notetaker} & \textbf{Question sets} \\
\midrule
S1 & & & & & \\
S2 & & & & & \\
\bottomrule
\end{tabularx}
\end{table}

% =====================================================================
\section{Analysis}
\label{sec:analysis}

\subsection{Empathy Maps}
\label{sec:empathy}

\begin{table}[H]
\centering
\caption{Empathy map: executives}
\label{tab:empathy-exec}
\small
\begin{tabularx}{\textwidth}{YY}
\toprule
\textbf{Says} & \textbf{Thinks} \\
\midrule
 & \\[3em]
\midrule
\textbf{Does} & \textbf{Feels} \\
\midrule
 & \\[3em]
\bottomrule
\end{tabularx}
\end{table}

\begin{table}[H]
\centering
\caption{Empathy map: team captains and experienced members}
\label{tab:empathy-captains}
\small
\begin{tabularx}{\textwidth}{YY}
\toprule
\textbf{Says} & \textbf{Thinks} \\
\midrule
 & \\[3em]
\midrule
\textbf{Does} & \textbf{Feels} \\
\midrule
 & \\[3em]
\bottomrule
\end{tabularx}
\end{table}

\subsection{Themes}
\label{sec:themes}

\subsection{Impact on Requirements}
\label{sec:impact}

\begin{table}[H]
\centering
\caption{Requirement changes resulting from the interviews}
\label{tab:impact}
\small
\begin{tabularx}{\textwidth}{Y L{2.4cm} L{2.6cm} L{2.0cm}}
\toprule
\textbf{Finding} & \textbf{Affected ID} & \textbf{Action} & \textbf{Issue} \\
\midrule
 & & & \\
 & & & \\
\bottomrule
\end{tabularx}
\end{table}

\subsection{Available Data for Testing}
\label{sec:test-data}
\todo{Past models and recorded failure loads identified in C4, and their use in
the VnV Plan.}

% =====================================================================
\section{Limitations}
\label{sec:limitations}
\todo{Sample size, selection, group-interview effects, interviewer bias.}

% =====================================================================
\newpage
\bibliographystyle{IEEEtran}
\bibliography{../../../refs/References}

\newpage
\appendix

\section{Raw Data: Interview Transcripts}
\label{app:transcripts}

\subsection{Session S1}
\label{app:s1}

\subsection{Session S2}
\label{app:s2}

\section{Use of Generative AI}
\label{app:ai}

\end{document}
